\documentclass[11pt]{article}
\usepackage[a4paper,margin=1in]{geometry}
\usepackage{amsmath,amssymb}
\usepackage{booktabs}
\usepackage{hyperref}
\hypersetup{hidelinks}

\newcommand{\wt}{\widetilde}
\newcommand{\Fcal}{\mathcal{F}}
\newcommand{\Scal}{\mathcal{S}}

\title{Corrected Eq. (29)--Eq. (41) Snippet for the Updated Shortcut Calculation}
\author{}
\date{7 June 2026}

\begin{document}
\maketitle

\section*{Scope}

This note records the corrected replacement formulae for the updated
shortcut-to-absorbing-target calculation.  The source checked is Luca's
2026-06-06 updated \texttt{Calculations.tex}.  The leading shortcut resolvent
sign has been corrected in the updated Eq. (29)--(31), but the correction has
not been propagated consistently into Eq. (32), Eq. (33), Eq. (40), Eq. (41),
and the final tail statement.

Let
\[
  \lambda=\beta(1-q).
\]
For a shortcut \(u\to v\) where \(v\) is the absorbing target, the shortcut
acts as killing in the transient state space: probability mass is removed from
the \(u\to u\) self-loop and absorbed at \(v\).

\section*{Corrected generating-function chain}

The updated Eq. (31) is the correct starting point:
\begin{equation}
  \wt S_{n_0}(n,z)
  =
  \wt W_{n_0}(n,z)
  -
  \frac{
    z\lambda \wt W_{n_0}(u,z)\wt W_u(n,z)
  }{
    1+z\lambda \wt W_u(u,z)
  }.
  \label{eq:snippet-S-correct}
\end{equation}

Summing over transient sites and using
\[
  \sum_n \wt W_r(n,z)=
  \frac{1-\wt{\Fcal}_r(v,z)}{1-z}
\]
gives the corrected replacement for Eq. (32):
\begin{equation}
  \boxed{
  \wt F_{n_0}(v,z)=
  \wt{\Fcal}_{n_0}(v,z)
  +
  \frac{
    z\beta(1-q)\wt W_{n_0}(u,z)
    \left[1-\wt{\Fcal}_u(v,z)\right]
  }{
    1+z\beta(1-q)\wt W_u(u,z)
  }}.
  \label{eq:snippet-F-correct}
\end{equation}

With
\begin{equation}
  \wt H_{N,|u-v|}(z)
  =
  \frac{\beta(1-q)}{q}
  \frac{1}{1+z\beta(1-q)\wt W_u(u,z)},
  \label{eq:snippet-H-def}
\end{equation}
the corrected replacement for Eq. (33) is
\begin{equation}
  \boxed{
  \wt F_{n_0}(v,z)=
  \wt{\Fcal}_{n_0}(v,z)
  +
  qz\,\wt W_{n_0}(u,z)
  \left[1-\wt{\Fcal}_u(v,z)\right]
  \wt H_{N,|u-v|}(z)}.
  \label{eq:snippet-F-H-correct}
\end{equation}

For \(N=2L\) and \(|u-v|=N/2\), the corrected \(H\) denominator is
\begin{equation}
  \wt H_{N,N/2}(z)
  =
  \frac{1}{\frac{q}{\beta(1-q)}+\phi(z)}.
  \label{eq:snippet-H-symmetric}
\end{equation}
This plus sign is consistent with \eqref{eq:snippet-H-def}.

\section*{Corrected time-domain convolution}

The corrected replacement for Eq. (40) is
\begin{equation}
  \boxed{
  F_{n_0}(v,t)=\Fcal_{n_0}(v,t)
  +q\sum_{t'=0}^{t-1}W_{n_0}(u,t-1-t')
  \sum_{t''=0}^{t'}
  \left[\delta_{t',t''}-\Fcal_u(v,t'-t'')\right]
  H_{N,N/2}(t'')}
  \label{eq:snippet-time-convolution}
\end{equation}
for \(t\geq1\), with \(F_{n_0}(v,0)=0\) when \(n_0\neq v\).

\section*{Clean replacement for Eq. (41)}

The long mixed \(\alpha,\gamma,s\) convolution form is not the safest final
form.  For the symmetric antipodal case, set \(N=2L\), \(v=0\), \(u=L\), and
let \(\rho\in\{0,\ldots,L\}\) be the folded distance from the initial site to
the target.  Define
\[
  a=\frac{q}{\beta(1-q)},\qquad
  y=1+\frac{1/z-1}{q},\qquad U_{-1}(y)=0.
\]
The corrected first-passage generating function reduces to
\begin{equation}
  \boxed{
  \wt F_\rho(z)
  =
  \frac{
    aT_{L-\rho}(y)+U_{\rho-1}(y)+U_{L-\rho-1}(y)
  }{
    aT_L(y)+U_{L-1}(y)
  }}.
  \label{eq:snippet-closed-form}
\end{equation}

Let
\[
  D(y)=aT_L(y)+U_{L-1}(y),\qquad
  N_\rho(y)=aT_{L-\rho}(y)+U_{\rho-1}(y)+U_{L-\rho-1}(y).
\]
If \(D(y_k)=0\) and \(s_k=1-q+qy_k\), then for simple roots
\begin{equation}
  \boxed{
  F_\rho(t)=\sum_k B_{\rho k}s_k^{t-1},
  \qquad
  B_{\rho k}=\frac{q\,N_\rho(y_k)}{D'(y_k)}
  }
  \label{eq:snippet-partial-fraction}
\end{equation}
for \(t\geq1\).  This is the recommended time-domain replacement for Eq. (41).

\section*{Corrected tail statement}

The apparent \(\alpha_\ell\) and \(\gamma_r\) poles in the long convolution
expansion are intermediate poles.  In the closed form
\eqref{eq:snippet-closed-form}, the ultimate tail is controlled by the largest
positive root \(s_1\) of \(D(y)=0\):
\begin{equation}
  \boxed{
  F_\rho(t)\sim B_{\rho1}s_1^{t-1},
  \qquad
  \frac{F_\rho(t+1)}{F_\rho(t)}\to s_1
  }.
  \label{eq:snippet-tail}
\end{equation}
Equivalently,
\begin{equation}
  \boxed{
  \tau_\beta^{-1}=-\log s_1,\qquad
  aT_L\left(1+\frac{s_1-1}{q}\right)
  +
  U_{L-1}\left(1+\frac{s_1-1}{q}\right)=0
  }.
  \label{eq:snippet-decay-rate}
\end{equation}

For weak shortcut strength,
\begin{equation}
  s_1
  =
  \alpha_1-\frac{2\beta(1-q)}{N}
  +O\!\left([\beta(1-q)]^2\right),
  \qquad
  1-s_1
  \sim
  \frac{q\pi^2}{2N^2}
  +
  \frac{2\beta(1-q)}{N}.
  \label{eq:snippet-weak-beta}
\end{equation}
Thus the generic \(t\alpha_1^t\) tail in the updated manuscript should be
removed.

\section*{Numerical checks}

For \(N=100\), \(q=2/3\), \(\beta=0.04\), shortcut \(u=6\to v=56\), and target
\(v=56\), direct finite-matrix checks give:
\begin{center}
\begin{tabular}{lc}
\toprule
Formula compared with the direct stochastic matrix & max absolute error \\
\midrule
correct plus/plus first-passage formula & \(2.08\times10^{-17}\) \\
updated Eq. (32) denominator sign & \(3.29\times10^{-3}\) \\
updated Eq. (33) minus prefactor & \(6.34\times10^{-2}\) \\
closed form \eqref{eq:snippet-closed-form} & \(7.63\times10^{-17}\) \\
\bottomrule
\end{tabular}
\end{center}

For the same \(N=100\), \(q=2/3\), the proposed equality
\[
  s_1=1-q+q\cos\frac{2\pi}{N}
\]
does not occur for finite positive \(\beta\) in the admissible range.  The
corrected dominant pole satisfies
\[
  \gamma_1<s_1<\alpha_1,
  \qquad
  \gamma_1=1-q+q\cos\frac{2\pi}{N},
  \qquad
  \alpha_1=1-q+q\cos\frac{\pi}{N}.
\]

Under the clear visual double-peak criterion used in the local scan, clear
double peaks occur up to about \(\beta=0.03\) for \(n_0=1,2,3\), and no clear
double peak remains at \(\beta_c=2q/[(1-q)N]=0.04\).

\end{document}
